% GENERATED FILE. Edit canonical lessons, then run npm run generate:books.
% canonical-source-hash: fnv1a64:b4cf7dd4a463e63c
% canonical-lessons: FR-C455-randonnee, FR-C455-randonneur, FR-C455-pente, FR-C455-gravir, FR-C455-trebucher, FR-R455-revision-448-455

\chapter{En randonnée}
\label{ch:fr-en-randonnee}
\hlchaptermodality{455}

\begin{chapteropening}
\textbf{By the end of this chapter:} \emph{I can talk about a hike: the hike, the hiker, the slope, climbing up and stumbling.}

Complete the last lesson of chapter 455: I can talk about a hike: the hike, the hiker, the slope, climbing up and stumbling.
\end{chapteropening}

\section[la randonnée]{la randonnée — the hike, the long walk}

\label{lesson:FR-C455-randonnee}



\begin{quote}
Before the new one: say the French for the mountain pass, then the French for the mountainside.
\end{quote}

\subsection*{You'll want to know: la randonnée}
\textbf{la randonnée} — \textquotedblleft{}the hike, the long walk\textquotedblright{}.

\begin{quote}
Le dimanche, nous faisons une randonnée en montagne.
\end{quote}

\begin{quote}
\emph{On Sundays we go for a hike in the mountains.}
\end{quote}

In speech, \emph{la rando}.

\begin{cousinweb}
\emph{randonnée} is built on \emph{randonner} 'to walk a long way' (EtymDB), from Old French \emph{randonner} 'to run headlong', usually traced to a Germanic word.
\end{cousinweb}

\begin{grammarlens}[title={What you've built}]
One more word for telling, explaining and describing at B1.
\end{grammarlens}

\subsection*{Your turn}
\begin{itemize}
\raggedright
  \item \emph{Say it:} \emph{la randonnée}
  \item \emph{Say it:} \emph{la randonnée}, once more
  \item \emph{Say it:} a sentence of your own with \emph{la randonnée}
  \item \emph{Recall:} say the French for the mountain pass, then the French for the mountainside, then say \emph{la randonnée} again
\end{itemize}

\begin{tcolorbox}[breakable,colback=teal!4,colframe=teal!35!black,title={Before you move on}]
What does la randonnée mean? (\textquotedblleft{}The hike, the long walk\textquotedblright{}.) Say it once more.
\end{tcolorbox}
\section[le randonneur]{le randonneur — the hiker}

\label{lesson:FR-C455-randonneur}



\begin{quote}
Before the new one: say the French for the mountainside, then the French for the hike.
\end{quote}

\subsection*{You'll want to know: le randonneur}
\textbf{le randonneur} — \textquotedblleft{}the hiker\textquotedblright{}.

\begin{quote}
Les randonneurs se reposent près du lac.
\end{quote}

\begin{quote}
\emph{The hikers are resting by the lake.}
\end{quote}

Feminine \emph{la randonneuse}.

\begin{cousinweb}
\emph{randonneur} is built on \emph{randonner} 'to hike' with \emph{-eur} (EtymDB).
\end{cousinweb}

\begin{grammarlens}[title={What you've built}]
One more word for telling, explaining and describing at B1.
\end{grammarlens}

\subsection*{Your turn}
\begin{itemize}
\raggedright
  \item \emph{Say it:} \emph{le randonneur}
  \item \emph{Say it:} \emph{le randonneur}, once more
  \item \emph{Say it:} a sentence of your own with \emph{le randonneur}
  \item \emph{Recall:} say the French for the mountainside, then the French for the hike, then say \emph{le randonneur} again
\end{itemize}

\begin{tcolorbox}[breakable,colback=teal!4,colframe=teal!35!black,title={Before you move on}]
What does le randonneur mean? (\textquotedblleft{}The hiker\textquotedblright{}.) Say it once more.
\end{tcolorbox}
\section[la pente]{la pente — the slope, the incline}

\label{lesson:FR-C455-pente}



\begin{quote}
Before the new one: say the French for the hike, then the French for the hiker.
\end{quote}

\subsection*{You'll want to know: la pente}
\textbf{la pente} — \textquotedblleft{}the slope, the incline\textquotedblright{}.

\begin{quote}
La pente est raide, alors nous montons doucement.
\end{quote}

\begin{quote}
\emph{The slope is steep, so we go up slowly.}
\end{quote}

\emph{En pente} — sloping.

\begin{cousinweb}
\emph{pente} goes back to Vulgar Latin \emph{pendita}, from Latin \emph{pendēre} 'to hang' (EtymDB).
\end{cousinweb}

\begin{grammarlens}[title={What you've built}]
One more word for telling, explaining and describing at B1.
\end{grammarlens}

\subsection*{Your turn}
\begin{itemize}
\raggedright
  \item \emph{Say it:} \emph{la pente}
  \item \emph{Say it:} \emph{la pente}, once more
  \item \emph{Say it:} a sentence of your own with \emph{la pente}
  \item \emph{Recall:} say the French for the hike, then the French for the hiker, then say \emph{la pente} again
\end{itemize}

\begin{tcolorbox}[breakable,colback=teal!4,colframe=teal!35!black,title={Before you move on}]
What does la pente mean? (\textquotedblleft{}The slope, the incline\textquotedblright{}.) Say it once more.
\end{tcolorbox}
\section[gravir]{gravir — to climb (up)}

\label{lesson:FR-C455-gravir}



\begin{quote}
Before the new one: say the French for the hiker, then the French for the slope.
\end{quote}

\subsection*{You'll want to know: gravir}
\textbf{gravir} — \textquotedblleft{}to climb (up)\textquotedblright{}.

\begin{quote}
Ils ont gravi la colline en une heure.
\end{quote}

\begin{quote}
\emph{They climbed the hill in an hour.}
\end{quote}

Conjugated like \emph{finir}. A bit formal.

\begin{cousinweb}
The origin of \emph{gravir} is not settled, and EtymDB does not list it.
\end{cousinweb}

\begin{grammarlens}[title={What you've built}]
One more word for telling, explaining and describing at B1.
\end{grammarlens}

\subsection*{Your turn}
\begin{itemize}
\raggedright
  \item \emph{Say it:} \emph{gravir}
  \item \emph{Say it:} \emph{gravir}, once more
  \item \emph{Say it:} a sentence of your own with \emph{gravir}
  \item \emph{Recall:} say the French for the hiker, then the French for the slope, then say \emph{gravir} again
\end{itemize}

\begin{tcolorbox}[breakable,colback=teal!4,colframe=teal!35!black,title={Before you move on}]
What does gravir mean? (\textquotedblleft{}To climb (up)\textquotedblright{}.) Say it once more.
\end{tcolorbox}
\section[trébucher]{trébucher — to stumble, to trip}

\label{lesson:FR-C455-trebucher}



\begin{quote}
Before the new one: say the French for the slope, then the French for to climb.
\end{quote}

\subsection*{You'll want to know: trébucher}
\textbf{trébucher} — \textquotedblleft{}to stumble, to trip\textquotedblright{}.

\begin{quote}
Sur le sentier, elle a trébuché sur une pierre.
\end{quote}

\begin{quote}
\emph{On the path she tripped over a stone.}
\end{quote}

\emph{Trébucher sur un mot} — to stumble over a word.

\begin{cousinweb}
\emph{trébucher} comes from Old French \emph{trebuchier} 'to topple' (EtymDB), from \emph{tres-}, from Latin \emph{trāns} 'across', and a Frankish word for 'trunk, body'.
\end{cousinweb}

\begin{grammarlens}[title={What you've built}]
One more word for telling, explaining and describing at B1.
\end{grammarlens}

\subsection*{Your turn}
\begin{itemize}
\raggedright
  \item \emph{Say it:} \emph{trébucher}
  \item \emph{Say it:} \emph{trébucher}, once more
  \item \emph{Say it:} a sentence of your own with \emph{trébucher}
  \item \emph{Recall:} say the French for the slope, then the French for to climb, then say \emph{trébucher} again
\end{itemize}

\begin{tcolorbox}[breakable,colback=teal!4,colframe=teal!35!black,title={Before you move on}]
What does trébucher mean? (\textquotedblleft{}To stumble, to trip\textquotedblright{}.) Say it once more.
\end{tcolorbox}
\section[(dialogue)]{Retour sur les chapitres 448–455}

\label{lesson:FR-R455-revision-448-455}



\begin{quote}
Say the French for the hairdresser and to shave.
\end{quote}

\subsection*{Your turn}
Say these aloud:

\begin{itemize}
\raggedright
  \item the hairdresser, to shave, the optician, the watchmaker, the jeweller
  \item the florist, the tobacconist, to put a stamp on, to send off, to tie up
  \item the cashier, to cash in, to charge, the card machine, cash
  \item the withdrawal, to credit, to settle, the overdraft, the statement
  \item the discount, the reduction, to haggle, to sell off cheap, the street sale
  \item the packaging, to wrap, the delivery man, the cardboard box, the registered letter
  \item the cliff, the plateau, the summit, the mountain pass, the mountainside
  \item the hike, the hiker, the slope, to climb, to stumble
\end{itemize}

\begin{tcolorbox}[breakable,colback=teal!4,colframe=teal!35!black,title={Before you move on}]
The hairdresser? (\textbf{le coiffeur}.) To shave? (\textbf{raser}.) The optician? (\textbf{l'opticien}.) The watchmaker? (\textbf{l'horloger}.) The jeweller? (\textbf{le bijoutier}.) The florist? (\textbf{le fleuriste}.) The tobacconist? (\textbf{le buraliste}.) To put a stamp on? (\textbf{affranchir}.) To send off? (\textbf{expédier}.) To tie up? (\textbf{ficeler}.) The cashier? (\textbf{le caissier}.) To cash in? (\textbf{encaisser}.) To charge? (\textbf{débiter}.) The card machine? (\textbf{le terminal}.) Cash? (\textbf{le liquide}.) The withdrawal? (\textbf{le retrait}.) To credit? (\textbf{créditer}.) To settle? (\textbf{régler}.) The overdraft? (\textbf{le découvert}.) The statement? (\textbf{le relevé}.) The discount? (\textbf{la remise}.) The reduction? (\textbf{le rabais}.) To haggle? (\textbf{marchander}.) To sell off cheap? (\textbf{brader}.) The street sale? (\textbf{la braderie}.) The packaging? (\textbf{l'emballage}.) To wrap? (\textbf{emballer}.) The delivery man? (\textbf{le livreur}.) The cardboard box? (\textbf{le carton}.) The registered letter? (\textbf{le recommandé}.) The cliff? (\textbf{la falaise}.) The plateau? (\textbf{le plateau}.) The summit? (\textbf{le sommet}.) The mountain pass? (\textbf{le col}.) The mountainside? (\textbf{le versant}.) The hike? (\textbf{la randonnée}.) The hiker? (\textbf{le randonneur}.) The slope? (\textbf{la pente}.) To climb? (\textbf{gravir}.) To stumble? (\textbf{trébucher}.)
\end{tcolorbox}
